% ============================================================
% APPENDICE B - TABELLA DEGLI ARMONICI
% ============================================================
\chapter{Tabella degli armonici per ogni nota}
\label{app:armonici}

La tabella dà i primi otto armonici di ciascuna nota. Il settimo armonico è calante rispetto al temperamento equabile (circa un sesto di tono): la nota scritta è la più vicina, con la grafia di settima minore della fondamentale. Le fondamentali alterate sono scritte con il nome più usato (Mi$\flat$ e non Re$\sharp$), perché gli armonici conservino la grafia corretta degli intervalli.

\begin{center}
\small
\renewcommand{\arraystretch}{1.15}
\begin{tabular}{lcccccccc}
\toprule
\textbf{Fondamentale} & \textbf{1} & \textbf{2} & \textbf{3} & \textbf{4} & \textbf{5} & \textbf{6} & \textbf{7} & \textbf{8} \\
\midrule
Do & Do & Do & Sol & Do & Mi & Sol & \Sibemolle{} & Do \\
\Dodiesis{} & \Dodiesis{} & \Dodiesis{} & \Soldiesis{} & \Dodiesis{} & \Midiesis{} & \Soldiesis{} & Si & \Dodiesis{} \\
Re & Re & Re & La & Re & \Fadiesis{} & La & Do & Re \\
\Mibemolle{} & \Mibemolle{} & \Mibemolle{} & \Sibemolle{} & \Mibemolle{} & Sol & \Sibemolle{} & \Rebemolle{} & \Mibemolle{} \\
Mi & Mi & Mi & Si & Mi & \Soldiesis{} & Si & Re & Mi \\
Fa & Fa & Fa & Do & Fa & La & Do & \Mibemolle{} & Fa \\
\Fadiesis{} & \Fadiesis{} & \Fadiesis{} & \Dodiesis{} & \Fadiesis{} & \Ladiesis{} & \Dodiesis{} & Mi & \Fadiesis{} \\
Sol & Sol & Sol & Re & Sol & Si & Re & Fa & Sol \\
\Labemolle{} & \Labemolle{} & \Labemolle{} & \Mibemolle{} & \Labemolle{} & Do & \Mibemolle{} & \Solbemolle{} & \Labemolle{} \\
La & La & La & Mi & La & \Dodiesis{} & Mi & Sol & La \\
\Sibemolle{} & \Sibemolle{} & \Sibemolle{} & Fa & \Sibemolle{} & Re & Fa & \Labemolle{} & \Sibemolle{} \\
Si & Si & Si & \Fadiesis{} & Si & \Rediesis{} & \Fadiesis{} & La & Si \\
\bottomrule
\end{tabular}
\end{center}
